\subsection{実験目的}
ここに実験の目的を記述する．提案手法の何を評価するのかを明確にする．

\subsection{実験条件}
ここに実験条件を記述する．データセット，評価指標，実験環境，比較対象などを整理する．

\subsection{実験結果}
ここに実験結果を記述する．表や図を用いて，結果を客観的に示す．

% ----- 図挿入の使用例（テンプレート参照） -----
% テンプレートファイル: templates/figure_template.tex
% 使い方の例: 以下のコメントを外すか、テンプレートを `\input{templates/figure_template.tex}` で読み込んでください。
%
% 単独図の直接例:
\begin{figure}[htbp]
  \centering
  \includegraphics[width=0.7\linewidth]{images/dummy.png}
  \caption{実験の代表的な結果}
  \label{fig:exp-result}
\end{figure}

% サブ図の例は templates/figure_template.tex を参照してください。


\subsection{結果の分析}
ここに実験結果から読み取れる傾向を記述する．考察章では解釈や原因の議論を中心に扱うため，本節では結果の整理に留める．
